\begin{circuitikz}[
        american,
        voltage dir=RP,
        >={Latex[length=5mm, width=2mm]},
        %>=latex,
        alt={Schematic diagram of a simple DC electrical circuit connected to a PLC output point labeled Output Point. On the right, a DC battery labeled V provides voltage to the circuit. Connected in series from the battery to the PLC interface on the top wire is an electrical load labeled Load. The PLC output module is enclosed in a dashed red box on the left, labeled NPN at the bottom, and contains an internal current source component. The return wire connects from the bottom of the PLC output module back to the negative terminal of the battery and is labeled Common.}
    ]

    \draw[
        BakerBlack
    ]
    (0,0)
    to[battery, color=BakerBlack, v_={$V$}] ++(0,2)
    to[R, color=BakerBlack, l_={Load}, font=\small, -o] ++(-3,0)
    coordinate(PLC_top)
    -- ++(-1,0)
    to[isource, color=BakerBlack] ++(0,-2)
    to[short, color=BakerBlack, -o] ++(1,0)
    coordinate(PLC_bottom)
    -- (0,0)
    node[
        pos=0.5,
        text=BakerBlack,
        font=\small,
        above
    ]{Common};

    \node[
        draw=BakerADARed,
        dashed,
        anchor=north east,
        minimum width=2cm,
        minimum height=3cm,
        label={above:{\small \textcolor{BakerBlack}{Output Point}}},
    ]
    (PLC)
    at ($(PLC_top)+(0,0.5)$)
    {};

    \node[
        text=BakerBlack,
        anchor=south
    ]
    at (PLC.south)
    {{\tiny NPN}};

\end{circuitikz}